\documentclass{article}
\usepackage{amsmath,amssymb,mathrsfs,mathtools}
\usepackage[margin=1in]{geometry}
\allowdisplaybreaks
\usepackage[UTF8]{ctex}
\begin{document}

\section*{自适应学习与路由策略优化}

请选择路由策略$\pi$，以最大化长期预期奖励率。该策略在论文中有所描述。
\begin{equation}
\frac{1}{t}\sum_{(ij)\in\mathcal L}\theta_{ij}\mathbb E(D_{ij}(t)).
\tag{1}
\end{equation}
在遵循以下排队动态和兼容关系的前提下，政策要求如下：
\begin{align}
\lim_{t\to\infty}\left(t\lambda_i-\sum_{j\in\mathcal S_i}\mathbb E^\theta_\pi(D_{ij}(t))\right)&<\infty,
&&\forall i\in\mathcal I,\tag{7}\\
t(\mu_j-\varepsilon)-\sum_{i\in\mathcal C_j}\mathbb E^\theta_\pi(D_{ij}(t))&\geq0,
&&\forall j\in\mathcal J,\quad\forall t\geq0.\tag{8}
\end{align}
该论文对$\varepsilon\geq0$以及任何收益向量都提出了这些要求。所选择的自适应构建方式使用固定的$\varepsilon>0$和伯努利收益，具体说明如下。路由策略决定了哪个客户由哪个服务器在何时提供服务。符号$\mathbb P^\theta_\pi,\mathbb E^\theta_\pi$表示在收益向量$\theta$下的概率分布和期望值。量$D_{ij}(t)$表示在时间$t\geq0$之前，在服务器$j$处完成类型$i$服务的次数总和。因此，累积的预期收益就是……
\begin{equation}
\sum_{(ij)\in\mathcal L}\theta_{ij}\mathbb E^\theta_\pi(D_{ij}(t)).
\tag{9}
\end{equation}
要求（8）指的是在服务器$j$上，预期的工作负载最多为$(\mu_j-\varepsilon)/\mu_j$。相关的要求（a）至（d）旨在遵循排队动态和兼容性关系，确保每个队列的稳定性，了解未知的或错误的$\theta_{ij}$，并有效处理探索与利用之间的平衡。

\paragraph{目标的适用范围。}
方程(1)保持不变，没有添加极限、liminf或折扣标准。来源将这一表达式称为“长期预期奖励率”，但没有使用任何限制运算符。在介绍(7)--(8)的句子中，来源将$\theta\in\mathbb R_{\geq0}$保留原样，同时将$\theta$称为向量；其在第2.1节至2.2节中的定义是$\theta\in\mathbb R_{\geq0}^{L}$。

\section*{固定的网络结构、随机输入以及收益限制}

固定的客户类型和服务器设置及其数量如下：
\[
\mathcal I=\{1,\ldots,I\},\qquad
\mathcal J=\{1,\ldots,J\},\qquad I,J\in\mathbb N_{>0}.
\]
打印出的兼容性线集及其大小的定义如下：
\[
\mathcal L=\{(ij):i\in\mathcal I,\ j\in\mathcal J\},\qquad
L=|\mathcal L|\in\mathbb N_{>0},\qquad L>I+J-1.
\]
兼容性邻域是
\begin{align}
\mathcal S_i&=\{j\in\mathcal J:(ij)\in\mathcal L\},\tag{2}\\
\mathcal C_j&=\{i\in\mathcal I:(ij)\in\mathcal L\}.\tag{3}
\end{align}
对于每一个$i\in\mathcal I$，有$i$个到达点，这些到达点构成了一个具有固定速率$\lambda_i>0$的泊松过程。如第4.3.1节所述，不同类型的客户之间的到达过程是独立的。客户在队列中等待$i$时间，直到被分配到一台兼容且空闲的服务器上。在时间$t=0$时，系统中没有客户。对于每一个$j\in\mathcal J$，服务器$j$的服务时间呈指数分布，其固定速率为$\mu_j>0$，且不同服务器之间以及不同客户之间的服务时间也是独立的。这些分布取决于服务器本身，而不是客户类型。
\[
\lambda=(\lambda_1,\ldots,\lambda_I),\qquad
\mu=(\mu_1,\ldots,\mu_J),\qquad
\mathcal S_A=\bigcup_{i\in A}\mathcal S_i.
\]
输入的稳定性假设是：
\begin{equation}
\sum_{i\in A}\lambda_i<\sum_{j\in\mathcal S_A}\mu_j,
\qquad\text{for any }\varnothing\ne A\subseteq\mathcal I.
\tag{4}
\end{equation}

\paragraph{印刷网络问题及选定的稳定性修正。}
$\mathcal L$的完整定义被保留下来，同时它也与其他内容存在冲突。原文中的(4)范围包括$A=\varnothing$，而需要$0<0$来补充。不过，实际的应用范围仅包括那一部分内容。这是一个带有注释的量化词修正，并非明确打印出的限定条件。源代码的第16页中，有$\lambda=(10,10)$、$\mu=(15,12)$和$\varepsilon=0.5$；所有这些非空切割都满足(4)的条件，而空切割则完全不符合这个条件。没有任何非空切割被放宽了。

在服务器$j$上完成$i$类型客户的服务后，会从$P_{\theta_{ij}}$中抽取金额$Y_{ij}$进行支付。
\[
\mathcal X\subseteq\mathbb N_{\geq0},\qquad
\theta_{ij}\in\mathbb R_{\geq0}\text{ unknown and finite},\qquad
\theta=(\theta_{ij})_{(ij)\in\mathcal L}\in\mathbb R_{\geq0}^{L}.
\]
所有行都使用相同的分布族，并且每行的均值也是相同的。第2.1节中的一般收益假设保持不变。
\begin{align}
I(\theta_{ij},\theta_{k\ell})
&:=\sum_{x\in\mathcal X}P_{\theta_{ij}}(x)
\ln\!\left(\frac{P_{\theta_{ij}}(x)}{P_{\theta_{k\ell}}(x)}\right),\tag{5}\\
0&<I(\theta_{ij},\theta_{k\ell})<\infty
\qquad\text{when }\theta_{ij}<\theta_{k\ell},
\quad (ij),(k\ell)\in\mathcal L.\tag{6}
\end{align}
对于所选的构造，第 4.1 节将收益函数进行了专门化处理。
\[
Y_{ij}\sim\operatorname{Ber}(\theta_{ij}),\qquad
\theta_{ij}\in[0,1],\qquad (ij)\in\mathcal L.
\]
线上的收益样本是独立且同分布的；正如第4.3.2节和附录A中所描述的，收益样本的抽取与系统的其他部分是相互独立的。

\paragraph{收益边界问题。}
封闭的伯努利参数区间被保留了下来。它并不自动满足之前的假设（6）：例如，将$\theta_{k\ell}=1$与$\theta_{ij}<1$进行比较时，会得到无限大的KL散度。没有添加任何开放区间的替换或端点约定。

\section*{LP 定义了动作空间，而不是策略目标。}

对于固定的$\theta\in\mathbb R_{\geq0}^{L}$和固定的$\varepsilon\geq0$，辅助静态问题为：
\begin{align}
\operatorname{LP}(\theta,\varepsilon):\quad
\max_x\quad&\sum_{(ij)\in\mathcal L}\theta_{ij}x_{ij},\tag{10a}\\
\text{subject to}\quad
&\sum_{j\in\mathcal S_i}x_{ij}=\lambda_i,
&&\forall i\in\mathcal I,\tag{10b}\\
&\sum_{i\in\mathcal C_j}x_{ij}\leq\mu_j-\varepsilon,
&&\forall j\in\mathcal J,\tag{10c}\\
&x_{ij}\geq0,
&&\forall(ij)\in\mathcal L.\tag{10d}
\end{align}
这里$x_{ij}\in[0,\lambda_i]$是一个路由速率变量，而$\varepsilon$是每个服务器上固定的松弛量。所选策略使用该问题的基本可行解作为行动方案，而不是用(10a)来替代(1)。

本文假设，该线性规划问题及其对偶问题的所有基本可行解都是非退化的。为了清晰起见，该假设中所指的对偶问题就是……
\begin{align}
\operatorname D(\theta,\varepsilon):\quad
\min_{v,w}\quad&\sum_{i\in\mathcal I}\lambda_i v_i+
\sum_{j\in\mathcal J}(\mu_j-\varepsilon)w_j,\tag{16a}\\
\text{subject to}\quad&v_i+w_j-\theta_{ij}\geq0,
&&\forall(ij)\in\mathcal L,\tag{16b}\\
&v_i\in\mathbb R,&&\forall i\in\mathcal I,\tag{16c}\\
&w_j\geq0,&&\forall j\in\mathcal J.\tag{16d}
\end{align}
这个双重问题只是非退化假设下所确定的辅助问题而已；其变量并不属于路由策略决策。来源方声称，在非退化假设下，最优的基本可行解是唯一的。

\subsection*{基础、森林、路由速率以及松弛的定义}

使用$\mathscr L\subseteq\mathcal L$表示选定的行列集合，使用$\mathscr J\subseteq\mathcal J$表示选定的服务器-松弛列集合。 $\mathscr L,\mathscr J$和$\mathcal L,\mathcal J$之间的区别是至关重要的。请将这些数学符号按原顺序原封不动地复制下来。同时，所有数字都需保持不变。 图$\mathcal G(\mathscr L,\mathscr J)$是一个 spanning forest，它由$(\mathcal I\cup\mathcal J,\mathcal L)$构成，当图中不存在循环时，该图包含$\mathcal I\cup\mathcal J$的每个顶点，并且其每个树都有一个唯一的根节点$j\in\mathscr J$。 父母和子女在客户类型和服务员之间交替出现。

以$v$为根节点的子树，包括$v$，用$\Cap(v)$来表示。
\[
\mathrm C(\Cap(v))=\Cap(v)\cap\mathcal I,\qquad
\mathrm S(\Cap(v))=\Cap(v)\cap\mathcal J.
\]
对于$\varepsilon\geq0$，引理1指出：
\[
B=\mathscr L\cup\mathscr J\text{ is a basis of }\operatorname{LP}(\theta,\varepsilon)
\quad\Longleftrightarrow\quad
\mathcal G(\mathscr L,\mathscr J)\text{ is a spanning forest}.
\]
附录 C.1 中提供了各列的定义。
\begin{align}
A'&\in\mathbb R^{I\times L},&
A'_{i,k\ell}&=\begin{cases}1,&i=k\text{ and }\ell\in\mathcal S_i,\\0,&\text{otherwise},\end{cases}\tag{112}\\
A''&\in\mathbb R^{J\times L},&
A''_{j,k\ell}&=\begin{cases}1,&j=\ell\text{ and }k\in\mathcal C_j,\\0,&\text{otherwise}.\end{cases}\tag{113}
\end{align}
用于定义这些基数的标准形式约束和松弛域如下：
\begin{equation}
\begin{bmatrix}A'&0\\A''&I\end{bmatrix}
\begin{bmatrix}x\\\sigma\end{bmatrix}
=\begin{bmatrix}\lambda\\\mu-\varepsilon\end{bmatrix},
\qquad x\in\mathbb R_{\geq0}^{L},\quad\sigma\in\mathbb R_{\geq0}^{J}.
\tag{111b}
\end{equation}
在这里，块矩阵中的$I$就是单位矩阵。对于基而言，$|B|=I+J$以及基矩阵具有满秩。
\begin{equation}
\mathbf B=
\begin{bmatrix}C_{\mathscr L}&D_{\mathscr J}\end{bmatrix}
=\begin{bmatrix}A'_{\mathscr L}&0_{\mathscr J}\\
A''_{\mathscr L}&I_{\mathscr J}\end{bmatrix}
\in\mathbb R^{(I+J)\times(I+J)}.
\tag{114}
\end{equation}
当其基本坐标为非负数时，基本解法是可行的；而当每个基本坐标都不为零时，该解法是有效的。

对于这样的基础，引理 2 明确给出了路由速率坐标$x\in\mathbb R^L$的顺序：
\begin{equation}
x_{ij}=\begin{cases}
\displaystyle\sum_{k\in\mathrm C(\Cap(i))}\lambda_k
-\sum_{\ell\in\mathrm S(\Cap(i))}(\mu_\ell-\varepsilon),
&\text{if }i\text{ is a child of }j\text{ in }\mathcal G(\mathscr L,\mathscr J),\\[6pt]
\displaystyle\sum_{\ell\in\mathrm S(\Cap(j))}(\mu_\ell-\varepsilon)
-\sum_{k\in\mathrm C(\Cap(j))}\lambda_k,
&\text{if }i\text{ is the parent of }j\text{ in }\mathcal G(\mathscr L,\mathscr J),\\[6pt]
0,&\text{if }(ij)\in\mathcal L\setminus\mathscr L.
\end{cases}
\tag{11}
\end{equation}
附录 C.2 中相关的松弛定义如下：
\[
\sigma_j:=\mu_j-\varepsilon-\sum_{i\in\mathcal C_j}x_{ij},
\qquad j\in\mathcal J.
\]
可行的基礎解決方案滿足了 (111b) 中的範圍要求。森林特徵描述和 (11) 是行動構建的定義，而非額外的路由策略。

\subsection*{允许使用的固定松弛量}

让$B=\mathscr L\cup\mathscr J$与$\mathscr L\subseteq\mathcal L$和$\mathscr J\subseteq\mathcal J$一起，构成$\operatorname{LP}(\theta,0)$的一个非退化可行基础。 引理 3 指出，如果……的话，那么$\operatorname{LP}(\theta,\varepsilon)$仍然是一个非退化的可行基础。
\begin{equation}
0\leq\varepsilon<\min_{j\in\mathcal J}
\frac{\displaystyle\sum_{\ell\in\mathrm S(\Cap(j))}\mu_\ell
-\displaystyle\sum_{k\in\mathrm S(\Cap(j))}\lambda_k}
{|\mathrm S(\Cap(j))|}.
\tag{12}
\end{equation}
第4.1节确定了网络$\mathcal I,\mathcal J,\mathcal L,\lambda,\mu$的位置，并选择了满足(12)条件的$\varepsilon>0$。引理3中的基于基础的前提条件被保留下来；没有对基础进行额外的限制。

\paragraph{印刷松弛问题。}
在(12)中，到达率的总和是由服务器集合$\mathrm S(\Cap(j))$来索引的，尽管$\lambda$是由客户类型来索引的。这是打印出来的表达式，也在(126)中重复出现；它并没有被改为$\mathrm C(\Cap(j))$。附录C.3中在(121)中打印了$\sigma_j^\varepsilon=\mu_j-\sum_{i\in\mathcal C_j}x_{ij}^\varepsilon$，但没有打印$-\varepsilon$；而在(127)中则打印了$\sigma_j^\varepsilon=\mu_j-\varepsilon-\sum_{i\in\mathcal C_j}x_{ij}^\varepsilon$。上述的操作松弛定义遵循的是(111b)和附录C.2中的规定，而不是与(121)相冲突的规定。

\section*{动作与事件驱动的学习与路由构建}

有限集合$\mathcal A$由(10)中列出的基本可行解组成。每个动作$a\in\mathcal A$都对应一个路由速率向量$x^a\in\mathbb R_{\geq0}^{L}$。这些向量取决于固定的网络、速率以及松弛量，而不是取决于收益信息的了解程度。定义如下：
\begin{align}
\mathcal L^a&:=\{(ij)\in\mathcal L:x^a_{ij}>0\},\tag{27}\\
r^a&:=\sum_{(ij)\in\mathcal L^a}x^a_{ij}\theta_{ij}\in\mathbb R_{\geq0},\tag{28}\\
d^a&:=\max_{\tilde a}r^{\tilde a}-r^a.\tag{29}
\end{align}
在(29)中，最大值指的是行动的数量。与最优基本可行解相对应的行动，被称为最优行动。

\subsection*{选定的回合长度规格}

请选择第 4.2 节中的理论性情节长度结构，而不是数值变体。其输入数量如下：
\[
x_{\min}:=\min_{a\in\mathcal A}\min_{(ij)\in\mathcal L^a}x^a_{ij},\qquad
\widetilde\lambda^a_j:=\sum_{i\in\mathcal C_j}x^a_{ij},
\qquad j\in\mathcal J,\quad a\in\mathcal A,
\]
\begin{equation}
\rho^a_j:=\widetilde\lambda^a_j/\mu_j\in(0,1),
\qquad j\in\mathcal J,\quad a\in\mathcal A.
\tag{34}
\end{equation}
请选择符合要求的固定值$\beta\in\mathbb R_{>1}$和$\alpha\in\mathbb R_{\geq1}$。
\begin{equation}
\alpha\geq\max_{j\in\mathcal J}\max_{a\in\mathcal A}
\left\{
\frac{3(\widetilde\lambda^a_j+\mu_j)\ln(\rho^a_j)
-2\sqrt{(\mu_j-\widetilde\lambda^a_j)^2+
9\widetilde\lambda^a_j\mu_j\ln^2(\rho^a_j)}}
{2(\mu_j-\widetilde\lambda^a_j)^2\ln(\rho^a_j)},\ 1
\right\}.
\tag{35}
\end{equation}
对于$H_0\in\mathbb R_{\geq1}$，预先确定的长度分别是：
\begin{equation}
H_k=\tau_k+H_0,\qquad
\tau_k=\alpha\ln^\beta(2Jk),\qquad
H_0\geq\max\left\{\frac4{x_{\min}},1\right\},
\qquad k=1,2,\ldots.
\tag{36}
\end{equation}
每集的时长均为$H_k\in\mathbb R_{>0}$。热身阶段之后是持续$H_0$的环节；路由算法在这两个阶段中都会运行。附录B中定义了累计的持续时间。
\begin{equation}
T_k:=\sum_{m=1}^{k}H_m,\qquad k\in\mathbb N_{\geq1}.
\tag{75}
\end{equation}
标量$T_k$与下面的行计数器$T_{ij}(k)$是不同的。

\paragraph{零负载边界问题。}
在(34)中，每个服务器和操作的积极负载间隔都被列出，但(10)允许使用未使用的服务器来进行操作，这种情况下使用$\widetilde\lambda^a_j=\rho^a_j=0$来表示。对于(35)中的对数计算，源代码中并未提供零负载的情况。间隔和公式都保持不变，没有限制条件，也没有删除未使用的服务器。

\subsection*{算法 1：自适应 UCB 队列路由}

初始化
\[
k=1,\qquad
T_{ij}(0)=0,\quad S_{ij}(0)=0,\quad\widehat\theta_{ij}(0)=0,\quad U_{ij}(0)=\infty
\quad\forall(ij)\in\mathcal L,
\qquad U^a(0)=\infty\quad\forall a\in\mathcal A.
\]
对于长度为$H_k$的回合$k=1,2,\ldots$：
\begin{enumerate}
\item 选择
\[
A_k=\operatorname*{argmax}_{a}U^a(k-1),
\]
任意打破连接。
\item 如果$k>1$和$A_k\ne A_{k-1}$，那么将每一个等待类型为$i$的顾客重新分配到服务器$j$的虚拟队列中，其概率如下：
\[
\frac{x^{A_k}_{ij}}{\lambda_i}.
\]
对于每一个$j\in\mathcal J$，按照顾客到达的时间顺序，将他们排列在那个虚拟队列中。同时，保留所有的数字。
\item 在整集的播放过程中，按照下面所指定的 FCFS RR$(A_k)$顺序来为顾客服务。打印出来的时间循环标题为“在时间$t=0,\ldots,H_k$”。
\item 对于每一个$(ij)\in\mathcal L$，统计第$k$集中所有实际完成的任务，包括在前一集中开始的任务。让$N^k_{ij}$负责统计这些任务，而$(Y^\ell_{ij})$则负责计算这些任务的报酬。保持原有的顺序不变。
\begin{align}
T_{ij}(k)&=T_{ij}(k-1)+N^k_{ij},\tag{30}\\
S_{ij}(k)&=S_{ij}(k-1)+\sum_{\ell=1}^{N^k_{ij}}Y^\ell_{ij},\\
\widehat\theta_{ij}(k)&=S_{ij}(k)/T_{ij}(k),\qquad T_{ij}(k)>0.\tag{31}
\end{align}
在零计数时，经验均值不会被计算或需要计算。其在$k=0$处的初始值会被保留下来；不会再施加新的零计数均值。对于每一行来说，
\begin{equation}
U_{ij}(k)=\begin{cases}
S_{ij}(k)/T_{ij}(k)+\sqrt{\ln(k)/T_{ij}(k)},&T_{ij}(k)>0,\\
\infty,&T_{ij}(k)=0.
\end{cases}\tag{32}
\end{equation}
\item 请计算每个动作的当前索引值：
\begin{equation}
U^a(k)=\sum_{(ij)\in\mathcal L^a}x^a_{ij}U_{ij}(k),\qquad a\in\mathcal A.\tag{33}
\end{equation}
只有正支持线可以进入这个总和，避免使用$0\cdot\infty$。
\end{enumerate}
这里的$N^k_{ij}$表示第$k$集中类型为$(ij)$的出发次数。论文中还将$T_{ij}(k)$定义为截至第$k\in\mathbb N_{\geq1}$集，在服务器$j$上发生的类型为$i$的出发总数；$\widehat\theta_{ij}(k)$则代表经验性收益均值；而$U_{ij}(k)$则代表其UCB估计值。所观察到的样本、计数器、估计值以及选定的行动都是随机策略量；而$\theta$则保持固定且未知。

\subsection*{算法 2：FCFS RR$(a)$}

对于固定动作$a\in\mathcal A$：
\begin{enumerate}
\item 每当有类型为$i$的顾客到达时，将该顾客分配到服务器$j$的虚拟队列中，其概率为……
\[
\frac{x^a_{ij}}{\lambda_i}.
\]
如果服务器处于空闲状态，则启动服务。
\item 在服务器$j$上，每次服务完成之后，都需要获取以下数据：
\[
Y_{ij}\sim\operatorname{Ber}(\theta_{ij}),
\]
其中，$i$是离开的顾客类型。如果服务器$j$的虚拟队列中有顾客，则从队列头部开始为顾客提供服务。
\end{enumerate}
虚拟队列采用先来先服务的原则。在任务变更时，算法1会重新分配等待的顾客，并恢复到达顺序；但它不指定已经处于服务中的顾客的重新分配。

\paragraph{原始冲突与选定的可执行源文件形式。}
算法1仅打印出所选支持线路和所选操作的更新信息，而它的(31)部分则缺少零计数的情况。 相比之下，第11页对$T$的定义中，使用了所有实际的数据，而附录C.10的第42页，方程(219)则使用了当前对$a$和$m$的估计值。 选择这些后源定义；它们并不与印刷版中仅用于选择的伪代码完全等价。 累计量 $S$ 是代数上的充分统计量。计数为正时，它与 $T$ 的比值等于式 (31)；计数为零时，式 (32) 直接给出无穷大。 在职的客户可以继续享受他们的服务；只有那些正在等待的客户才会被重新分配。 后续完成将基于实际线路进行。 没有$A_0$被分配：第一集的队列为空，因此省略空重新分配对路由没有影响。 通过论文的连续时间动力学和实际持续时间，解释这种看似离散的时间循环。
所有原始陈述都保留在候选变更日志中。

\section*{链接的累积奖励基准}

对于满足(12)条件的固定$\varepsilon>0$，令$x^\theta\in\mathbb R_{\geq0}^{L}$为(10)式的最优基本可行解。该来源定义了在时间$t\geq0$时，路由策略的预期遗憾值。
\begin{equation}
R^\theta_\pi(t):=
\sum_{(ij)\in\mathcal L}\theta_{ij}
\left(tx^\theta_{ij}-\mathbb E^\theta_\pi(D_{ij}(t))\right).
\tag{13}
\end{equation}
这是与保留的自适应策略相关的基准，而非替代目标或额外的可行性约束。遗憾分析子问题以及数值基准策略并不属于这一模型的一部分。

\paragraph{来源范围。}
该公式的推导遵循了所提供的原始论文，具体位于物理页面2、5-8、10-11，以及附录B、C.1-C.3和C.10中的定义材料（第42页，方程(219)）。文中引用的外部LP和森林结果[6]和[37]并未经过外部来源的验证；论文中给出的定义公式则如上所述。

\end{document}
